\resume
\selectlanguage{french}
\vspace*{-6cm}
\begin{abstract}

Dans un contexte où l'éducation numérique transforme les méthodes pédagogiques traditionnelles, ce mémoire présente la conception et le développement d'une application mobile innovante de gestion d'événements interactifs éducatifs. L'objectif principal est de créer une solution technologique qui améliore significativement l'engagement des participants et l'efficacité pédagogique lors de formations, conférences et cours interactifs.

Notre approche méthodologique s'articule autour de l'analyse comparative des solutions existantes (Kahoot, Slido, Poll For All), de l'identification de leurs limites, et du développement d'une architecture robuste utilisant Flutter pour l'interface mobile et Supabase comme Backend-as-a-Service (BaaS). L'application développée propose une solution complète avec 16 interfaces utilisateur organisées en cinq catégories fonctionnelles : authentification et gestion de compte, gestion des événements, modules interactifs (quiz et sondages), participation en temps réel, et administration des participants.

Les résultats obtenus démontrent l'efficacité de cette approche technologique pour créer des événements véritablement interactifs. L'intégration de Supabase permet des fonctionnalités temps réel (synchronisation instantanée des résultats, chat interactif, notifications de participation) avec une base de données PostgreSQL performante et un système d'authentification sécurisé. L'application offre aux organisateurs un tableau de bord complet pour suivre l'engagement et analyser les performances, tout en garantissant aux participants une expérience utilisateur fluide et engageante.

Cette recherche contribue à l'avancement des connaissances sur les applications mobiles éducatives interactives et démontre le potentiel des technologies BaaS pour accélérer le développement d'applications pédagogiques modernes. Elle ouvre des perspectives prometteuses pour la transformation numérique de l'éducation en Afrique et au-delà.

\vspace{0.5cm}

\textbf{Mots clés} : Application mobile éducative, Événements interactifs, Supabase, Backend-as-a-Service, Flutter, Temps réel, Quiz pédagogiques, Sondages, Engagement des participants

\end{abstract}

\newpage
\thispagestyle{empty}
\selectlanguage{english}
\addcontentsline{toc}{chapter}{Abstract}
\begin{abstract}

In a context where digital education is transforming traditional pedagogical methods, this thesis presents the design and development of an innovative mobile application for managing interactive educational events. The main objective is to create a technological solution that significantly improves participant engagement and pedagogical effectiveness during training sessions, conferences, and interactive courses.

Our methodological approach is structured around comparative analysis of existing solutions (Kahoot, Slido, Poll For All), identification of their limitations, and development of a robust architecture using Flutter for the mobile interface and Supabase as Backend-as-a-Service (BaaS). The developed application provides a comprehensive solution with 16 user interfaces organized into five functional categories: authentication and account management, event management, interactive modules (quizzes and polls), real-time participation, and participant administration.

The results obtained demonstrate the effectiveness of this technological approach for creating truly interactive events. Supabase integration enables real-time functionalities (instant result synchronization, interactive chat, participation notifications) with a high-performance PostgreSQL database and secure authentication system. The application provides organizers with a comprehensive dashboard to track engagement and analyze performance, while ensuring participants have a smooth and engaging user experience.

This research contributes to advancing knowledge on interactive educational mobile applications and demonstrates the potential of BaaS technologies to accelerate the development of modern pedagogical applications. It opens promising perspectives for the digital transformation of education in Africa and beyond.

\vspace{0.5cm}

\textbf{Key words} : Educational mobile application, Interactive events, Supabase, Backend-as-a-Service, Flutter, Real-time, Educational quizzes, Polls, Participant engagement

\end{abstract}
